\section{Capítulo~\ref{sec:dofa}. Aprendizaje con \nt{DOPA}}

Los mecanismos de la sinapsis (porciones, detector de coincidencias, calcio, ventanas de plasticidad) y el comportamiento de la dopamina como error de predicción están medidos directamente; que las reglas lleven al gradiente y lo que cuesta la difusión son teoremas; el resultado del aprendizaje de elección es un cálculo con el modelo del libro; el circuito que calcula el error es una hipótesis.

{\footnotesize
\setlength{\tabcolsep}{3pt}\renewcommand{\arraystretch}{1.15}
\begin{longtable}{>{\raggedright}p{0.5cm}>{\raggedright}p{4.0cm}>{\raggedright}p{5.6cm}>{\raggedright}p{2.3cm}>{\raggedright\arraybackslash}p{1.7cm}}
\toprule
N.º & Afirmación & Experimento y qué se midió & Fuente & Estado \\
\midrule\endfirsthead
\toprule
N.º & Afirmación & Experimento y qué se midió & Fuente & Estado \\
\midrule\endhead
1 & En el cerebro no se ha hallado retropropagación del error en la forma de las redes artificiales & No hay experimento directo: es una afirmación de ausencia. Las revisiones analizan qué exige el algoritmo (conexiones de vuelta que repitan las de ida, una fase aparte) y qué aproximaciones biológicas se han propuesto. & \cite{Lillicrap2020, WhittingtonBogacz2019} & hipótesis \\
2 & El peso se descompone en número de sitios de liberación, probabilidad de liberación y respuesta a una porción: $w=N_s\,p\,A_1$~\eqref{eq:quantal} & Unión neuromuscular de rana con liberación reducida: la respuesta del destinatario varía en escalones múltiplos de la respuesta miniatura espontánea a una porción; el número de porciones por espiga sigue la ley de Poisson. & \cite{delCastillo1954} & medido \\
3 & El receptor del destinatario es un detector de coincidencias; el calcio dispara el cambio de peso & Patch-clamp en neuronas de ratón: los iones Mg$^{2+}$ taponan el canal abierto por el glutamato en reposo y salen con la despolarización. Rebanadas de hipocampo: un quelante de calcio en el destinatario bloquea la potenciación a largo plazo, y liberar calcio con luz dentro de la célula refuerza por sí solo la transmisión. & \cite{Nowak1984, Malenka1988calcium} & medido \\
4 & La decisión la ejecuta cualquiera de los lados: crece $A_1$ en el destinatario, cambia $p$ en el emisor & El glutamato liberado con luz sobre una sola espina provoca un crecimiento rápido y duradero de la espina y de la corriente por sus receptores; la potenciación va acompañada de la inserción de receptores AMPA. La despolarización del destinatario suprime temporalmente la liberación en el emisor; el efecto lo transmiten endocannabinoides que cruzan la hendidura hacia atrás (mostrado en entradas \nt{GABA}). La depresión a corto plazo depende de la probabilidad de liberación en el emisor. & \cite{Matsuzaki2004, Malinow2002, WilsonNicoll2001, Tsodyks1997} & medido \\
5 & La sinapsis se refuerza cuando el emisor y el destinatario están activos a la vez~\eqref{eq:hebb} & Conejo anestesiado: tras series breves de espigas por la vía de entrada del hipocampo, la respuesta queda reforzada de $30$~min a $10$~h. En una rebanada de hipocampo, las espigas del destinatario refuerzan solo las sinapsis activas en ese mismo momento. & \cite{Bliss1973LTP, Kelso1986hebb} & medido \\
6 & La regla~\eqref{eq:oja} halla el vector propio principal de las correlaciones de las entradas (proposición~\ref{prop:oja}) & Teorema; demostración en el capítulo. No hay experimento en que una neurona haya aprendido la componente principal de sus entradas; el mecanismo que limita el crecimiento de los pesos en la sinapsis no se conoce. & \cite{Oja1982} & teoría \\
7 & Hebb temporal: ventana de unos $\pm20\ms$ (fig.~\ref{fig:stdp}); de secuencias repetidas crecen cadenas & Pares de neuronas de hipocampo en cultivo: una espiga del destinatario hasta $20\ms$ después de la del emisor da un refuerzo duradero; hasta $20\ms$ antes, un debilitamiento; ambos dependen de los receptores NMDA. La relación $A_-=0.6A_+$ de la figura es ilustrativa. Que así crezcan cadenas en la red no se ha medido directamente. & \cite{BiPoo1998} & medido \\
8 & Traza~\eqref{eq:threefactor}: la dopamina que llega $1$--$2\,\text{s}$ después de la actividad conjunta refuerza la conexión; más tarde, no (proposición~\ref{prop:threefactor}) & Rebanadas de estriado, estimulación optogenética separada de las entradas glutamatérgicas y dopaminérgicas: la espina crece solo si la dopamina llega $0.3$--$2\,\text{s}$ después de la entrada glutamatérgica; la ventana la fija la subida y caída rápidas del AMPc. En la corteza, el apareamiento deja trazas separadas para el refuerzo y el debilitamiento, que los receptores de monoaminas convierten en cambios duraderos en una ventana del orden de segundos. & \cite{Yagishita2014, He2015eligibility} & medido \\
9 & También hay ventanas de segundos sin dopamina: la tercera señal es una excitación fuerte del destinatario & Ratón vivo, campo CA1: un solo evento de meseta en el destinatario refuerza las entradas llegadas segundos antes y después, y crea un campo de lugar en una sola pasada. & \cite{Bittner2017} & medido \\
10 & El paso medio de la regla de tres factores sigue el gradiente de la recompensa~\eqref{eq:perturbation} (proposición~\ref{prop:gradient}) & Teorema del gradiente para el aprendizaje por perturbaciones aleatorias; en el capítulo se deduce para ruido pequeño. & \cite{Williams1992} & teoría \\
11 & El ruido es búsqueda: la red aprende de sus propias desviaciones aleatorias & Diamantes mandarín jóvenes: desactivar el núcleo de salida del circuito de los ganglios basales elimina de forma brusca y reversible la variabilidad del canto. Diamantes adultos: se aplica una perturbación a las ejecuciones de una sílaba cuyo tono queda a un lado del medio, y las aves desplazan pronto el tono hacia el otro lado: aprenden de su propia dispersión. & \cite{Olveczky2005, TumerBrainard2007} & medido \\
12 & La elección entre dos se aprende con $\tau_e=1\,\text{s}$ y $\delta=R-\bar R$; no se aprende con $\tau_e=50\ms$; sin restar, los pesos crecen sin límite, y con una tasa de aprendizaje alta la red puede fijar la peor elección & \texttt{models/dofa\_learning.py}, $\eta=0.05$, $500$ intentos, $6$ semillas. $\tau_e=1\,\text{s}$: fracción de elecciones de $A$ $0.65$--$0.79\to0.96$--$1.00$, pesos $0.85$--$1.33$. $\tau_e=50\ms$: $0.38$--$0.55$, los pesos no cambian. $\delta=R$: peso del ganador $860$--$1190$. $\delta=R$, $\eta=0.5$, $3$ semillas: una se fijó en $B$ ($0.04\to0$). El script imprime los cuatro casos; el recálculo coincidió con el capítulo. & \texttt{models/\allowbreak dofa\_\allowbreak learning.py} & modelo \\
13 & El tiempo de aprendizaje con una única señal común crece en proporción al número de neuronas ruidosas (proposición~\ref{prop:broadcastcost}) & Curvas de aprendizaje de una red lineal: la perturbación de neuronas es más lenta que el gradiente exacto tantas veces como neuronas de salida hay, y la perturbación de pesos, además, tantas veces como entradas. & \cite{Werfel2005} & teoría \\
14 & Las salidas de \nt{DOPA} van sobre todo a las regiones de elección de acciones & Reconstrucción completa de los axones de neuronas \nt{DOPA} nigroestriatales individuales de rata: las ramas de una célula cubren del $0.45$ al $5.7\%$ del volumen del estriado (en promedio, el $2.7\%$); fuera del estriado casi no hay ramas. & \cite{Matsuda2009} & medido \\
15 & La corteza aprende a predecir su entrada, y el error se calcula en el sitio & Revisión: en la corteza sensorial hay respuestas al desajuste entre la entrada esperada y la recibida. Que sea una regla general de aprendizaje de la corteza no está demostrado. & \cite{Keller2018} & hipótesis \\
16 & La dopamina se comporta como el error~\eqref{eq:ctd}: pico, traslado a la señal previa, valle (fig.~\ref{fig:ctdcases}) & Neuronas \nt{DOPA} de mono: pico ante un jugo inesperado; tras el aprendizaje, pico ante la señal previa y ninguna respuesta al jugo prometido; valle cuando el jugo no llega. La forma continua~\eqref{eq:ctd} es teoría. & \cite{Schultz1997, Doya2000} & medido \\
17 & Circuito que calcula $\dd V/\dd t$ y resta la expectativa & Ratones, registro y optogenética en el área tegmental ventral: las neuronas \nt{DOPA} restan la expectativa de la recompensa; las neuronas \nt{GABA} vecinas las inhiben cuando se espera la recompensa, y excitarlas o inhibirlas cambia el error. Cómo se calcula la derivada no se ha mostrado. & \cite{Eshel2015arithmetic} & hipótesis \\
18 & La neurona \nt{DOPA} es un marcapasos; los valles son menos profundos que los picos & En la rebanada, las neuronas \nt{DOPA} descargan de forma espontánea y muy regular, sobre un potencial marcapasos lento. En el mono, la frecuencia sigue cuantitativamente el error positivo, y el negativo lo transmite solo débilmente. & \cite{GraceOnn1989, Bayer2005} & medido \\
19 & Actor y crítico (fig.~\ref{fig:actorcritic}) & El algoritmo procede de la teoría del aprendizaje por refuerzo. En humanos (RMf), el estriado ventral se comporta como crítico con elección y sin ella; el dorsal, solo cuando hay que elegir acciones. & \cite{SuttonBarto2018, ODoherty2004actor} & hipótesis \\
20 & El signo de $\delta$ en la regla está invertido en las neuronas con la segunda familia de receptores & Rebanadas de estriado: en las neuronas con receptores D1 el apareamiento da un refuerzo que requiere D1; en las neuronas con D2, un debilitamiento que requiere D2. & \cite{Shen2008} & medido \\
\bottomrule
\end{longtable}}
